\chapter{Introduction}

\noindent This document describes FPAP, the Framebuffer Abstraction API. How
it must be implemented and how it may be used. It is assumed the reader has
a base familiarity with the C programming language.

FPAP is a series of functionw which aim to abstract startup issues when it
comes to drawing pixels directly to the screen to a small set of functions
which may be quickly and easily used to setup a render loop for new projects.

This means an implementer is most concerned with the specific method of
acquiring (or simulating) a framebuffer from the host and then applying the
various transforms from user input to hardware framebuffer.

\section{Versioning}

\noindent A version of this specification is defined in the form:
``V<major>.<minor>.<patch>''. Where each value is a two digit hexadecimal
value. An increment in the major version specifies that there has been a
change which breaks backwards compatibility (removing function behaviour,
changing identifier names, changing specification definitions, adding/removing
undefined behaviour). Whereas an increment in the minor version specifies an
addition in a way which does not break backwards compatibility, while a patch
specifies a small change to spelling, wording, or functionality that should
not effect existing implementations.

\section{Definitions}

\definitions
\name{access}
During execution time: To read or modify the value of an object.
\next{address}
A numeric value designating an ordered byte of memory.
\next{addressable}
(of an object): In a region of memory able to be represented with an address.
\next{arguments}
The ordered list of objects passed to a function during exececution.
\next{behaviour}
External behaviour or action.
\next{implementation defined behaviour}
Behaviour which is unspecified by this standard (within bounds) and documented
by the implementation.
\next{undefined behaviour}
Behaviour which is left entirely to the implementation, including all behaviour
not written in this here document. This may or may not be documented by the
implementation.
\next{bit}
Unit of storage which may store an object of two unique values (need not be
addressable).
\next{byte}
An addressable unit of storage which may store an object with at least 128
unique values, the specific number of bits is left unspecified.
\next{character}
An element of the American Standard Code for Information Interchange (ASCII)
as specified by ANSI INCITS 4-1986 (R2002). This is specified by a single byte
value.
\next{function}
An interface between the user and the framebuffer, window, and window system
which is passed an ordered set of objects with defined types (its arguments)
and passes back a value of predetermined type (its return value).
\next{implementation}
A particular software on a particular hardware for which this standard is
interpretted.
\next{implementation limit}
A restriction imposed by the implementation onto the program.
\next{object}
A region of storage during execution the contents of which may represent
values.
\next{prohibition}
A statement, behaviour, value, function or otherwise, which this standard
requires the implementation to handle in a particular manner described with
the qualifier `shall not'.
\next{requirement}
A statement, behaviour, value, function or otherwise, which this standard
requires the implementation to handle in a particular manner described with
the qualifier `shall'.
\next{type}
A qualifier denoting the method of accessing an object, including its size and
the kinds of operands which are defined.
\next{value}
The precise interpretation of an object when interpretted as a specific type.
\next{implementation defined value}
A value which is unspecified by this standard (within bounds) and documented by
the implementation.
\next{undefined value}
A value which is left entirely to the implementation. This may or may not be
documented by the implementation.
\next{window}
The system by which the region generated by the window system is accessed. It
need not exist, however an implementation shall act (to the user) as though it
does.
\next{window system}
The system by which fpap generates a region of the screen to assign to a
framebuffer. This system need not exist, however an implementation shall act
(to the user) as though it does.
\edefinitions
\vfill\break
